\documentclass{article}

\usepackage[utf8]{inputenc}
\usepackage[T1]{fontenc}
\usepackage{lmodern}
\usepackage[polish]{babel}
\usepackage[margin=1in]{geometry}

\usepackage{amsmath}
\usepackage{amsfonts}
\usepackage{amssymb}

\usepackage{booktabs} % Lepsza jakość tabel
\usepackage{graphicx}
\usepackage{float}
\usepackage{longtable}  % Dla tabel wielostronicowych

\usepackage{algorithm}
\usepackage{algpseudocode}
\usepackage{amsmath}
\usepackage{siunitx} 
\usepackage{algorithm}
\usepackage{algpseudocode}
\newcommand{\To}{\textbf{to }}
\newcommand{\DownTo}{\textbf{down to }}
\usepackage{amsmath}
\sisetup{
    output-exponent-marker = e,
    bracket-numbers = false,
    group-separator = {\,}, 
    scientific-notation = true
}

% Inne
\usepackage{hyperref} 
\hypersetup{
    colorlinks=true,
    linkcolor=blue,
    filecolor=magenta,      
    urlcolor=cyan,
    pdftitle={Obliczenia Naukowe - Lista 1},
    pdfauthor={Wojciech Typer},
}

\title{Sprawozdanie z Laboratorium\\Obliczenia Naukowe - Lista 5}
\author{Wojciech Typer}

\begin{document}
\maketitle
\section*{Zadanie 1}
\subsection*{Cel zadania}
Celem zadania jest napisanie funkcji rozwiązującej układ $ax = b$ metodą eliminacji Gaussa, 
uwzgledniając specyficzną postać macierzy $A$, która jest macierzą rzadką:
\[
A = 
\begin{pmatrix}
A_1    & C_1    & 0      & 0      & 0      & \dots  & 0      \\
B_2    & A_2    & C_2    & 0      & 0      & \dots  & 0      \\
0      & B_3    & A_3    & C_3    & 0      & \dots  & 0      \\
\vdots & \ddots & \ddots & \ddots & \ddots & \ddots & \vdots \\
0      & \dots  & 0      & B_{v-2} & A_{v-2} & C_{v-2} & 0      \\
0      & \dots  & 0      & 0      & B_{v-1} & A_{v-1} & C_{v-1} \\
0      & \dots  & 0      & 0      & 0      & B_v     & A_v
\end{pmatrix},
\]
dla dwóch wariantów:
\begin{itemize}
    \item bez wyboru elementu głównego
    \item z częściowym wyborem elementu głównego
\end{itemize}

\subsection*{Struktura danych i złożoność obliczeniowa} Macierz A jest macierzą rzadką o specyficznej strukturze blokowej, 
więc wystarczy tylko przechowywać jej niezerowe elementy. W programie wykorzystano strukturę słownika, co pozwala na alokację pamięci 
wyłącznie dla elementów niezerowych, redukując złożoność pamięciową do $O(n \cdot l)$. Złożoność czasowa procesów eliminacji i dekompozycji 
wynosi $O(n \cdot l^2)$, co dla małych, stałych wartości l jest złożonością liniową O(n). Wynika to z faktu, że pętle algorytmu ograniczają się 
jedynie do niezerowego pasma macierzy, pomijając operacje na pozostałych elementach. Pozwala to na stabilne i szybkie rozwiązanie układu nawet dla 
bardzo dużych wartości n.

\newpage
\subsection*{Implementacja}

\subsubsection*{Idea algorytmu eliminacji Gaussa bez wyboru elementu głównego}
Algorytm to „inteligentna” wersja eliminacji Gaussa, 
która wykorzystuje fakt, że macierz ma wartości inne niż zero tylko blisko swojej głównej przekątnej. 
Zamiast przeliczać całe wiersze i kolumny, program skupia się wyłącznie na wąskim pasmie o szerokości l, całkowicie ignorując ogromne obszary 
wypełnione zerami. Dzięki temu unikamy wykonywania milionów zbędnych działań (mnożenia przez zero), co pozwala błyskawicznie rozwiązać układ równań. 
Na samym końcu wynik x wyliczany jest od tyłu, również przy wykorzystaniu jedynie tych elementów, które znajdują się w pasmie.

\begin{algorithm}
\caption{Eliminacja Gaussa dla macierzy blokowej}
\begin{algorithmic}[1]
\Procedure{SolveGauss}{$n, l, A, b$}
    \For{$k \gets 1$ \textbf{ to } $n-1$}
        \State $last\_row \gets \min(k + l, n)$
        
        \For{$i \gets k+1$ \textbf{ to } $last\_row$}
            \If{$|A_{k,k}| < \epsilon$}
                \State \Return \textbf{error}("Zero na przekątnej")
            \EndIf
            
            \State $factor \gets A_{i,k} / A_{k,k}$
            \State $A_{i,k} \gets 0$
            
            \For{$j \gets k+1$ \textbf{ to } $\min(k + l, n)$}
                \State $A_{i,j} \gets A_{i,j} - factor \cdot A_{k,j}$
            \EndFor
            \State $b_i \gets b_i - factor \cdot b_k$
        \EndFor
    \EndFor

    \State $x \gets \text{new vector of size } n$
    \For{$i \gets n$ \textbf{ down to } $1$}
        \State $sum\_val \gets 0$
        \For{$j \gets i+1$ \textbf{ to } $\min(i + l, n)$}
            \State $sum\_val \gets sum\_val + A_{i,j} \cdot x_j$
        \EndFor
        \State $x_i \gets (b_i - sum\_val) / A_{i,i}$
    \EndFor
    \State \Return $x$
\EndProcedure
\end{algorithmic}
\end{algorithm}

\newpage

\subsubsection*{Idea eliminacji Gaussa z częściowym wyborem elementu głównego} Ta wersja algorytmu jest bezpieczniejsza dla wyników, 
ponieważ w każdym kroku szuka największej możliwej liczby w aktualnej kolumnie (tzw. elementu głównego) i zamienia wiersze tak, aby to ona 
znalazła się na przekątnej. Dzięki temu unikamy dzielenia przez bardzo małe liczby, co drastycznie poprawia dokładność obliczeń i chroni przed 
błędami zaokrągleń.

\begin{algorithm}
\caption{Eliminacja Gaussa z częściowym wyborem elementu głównego}
\begin{algorithmic}[1]
\Procedure{SolveGaussPivot}{$n, l, A, b$}
    \State $p \gets [1, 2, \dots, n]$ \Comment{Inicjalizacja wektora permutacji}
    
    \For{$k \gets 1$ \textbf{ to } $n-1$}
        \State $max\_val \gets 0.0$
        \State $pivot\_row \gets k$
        
        \Comment{Wybór elementu głównego w obrębie pasma $l$}
        \For{$i \gets k$ \textbf{ to } $\min(k + l, n)$}
            \If{$|A_{p_i, k}| > max\_val$}
                \State $max\_val \gets |A_{p_i, k}|$
                \State $pivot\_row \gets i$
            \EndIf
        \EndFor
        
        \State $\text{swap}(p_k, p_{pivot\_row})$ \Comment{Zamiana wierszy w wektorze permutacji}
        
        \For{$i \gets k+1$ \textbf{ to } $\min(k + l, n)$}
            \State $factor \gets A_{p_i, k} / A_{p_k, k}$
            \State $A_{p_i, k} \gets 0.0$
            
            \Comment{Aktualizacja wiersza - pasmo może wzrosnąć do $2l$}
            \For{$j \gets k+1$ \textbf{ to } $\min(k + 2l, n)$}
                \State $A_{p_i, j} \gets A_{p_i, j} - factor \cdot A_{p_k, j}$
            \EndFor
            \State $b_{p_i} \gets b_{p_i} - factor \cdot b_{p_k}$
        \EndFor
    \EndFor

    \State $x \gets \text{nowy wektor rozmiaru } n$
    \Comment{Podstawienie wsteczne z uwzględnieniem permutacji i pasma $2l$}
    \For{$i \gets n$ \textbf{ down to } $1$}
        \State $sum\_val \gets 0.0$
        \For{$j \gets i+1$ \textbf{ to } $\min(i + 2l, n)$}
            \State $sum\_val \gets sum\_val + A_{p_i, j} \cdot x_j$
        \EndFor
        \State $x_i \gets (b_{p_i} - sum\_val) / A_{p_i, i}$
    \EndFor
    
    \State \Return $x$
\EndProcedure
\end{algorithmic}
\end{algorithm}


\end{document}